\documentclass[11pt,a4paper]{article}
\usepackage[T1]{fontenc}
\usepackage{lmodern,microtype,amsmath,amssymb,amsthm,booktabs}
\usepackage[margin=28mm]{geometry}
\usepackage[hidelinks]{hyperref}
\newcommand{\E}{\mathbb E}
\newcommand{\Pp}{\mathbb P}
\newcommand{\CE}{C_{\mathrm E}}
\newtheorem{proposition}{Proposition}
\newtheorem{lemma}[proposition]{Lemma}
\title{Count-specific quartic majorants on a continuous parameter cell}
\author{Research continuation accompanying journal-v21}
\date{5 November 2026}
\begin{document}
\maketitle
\begin{abstract}
We give a finite exact certificate criterion for a conditional Berry--Esseen
comparison throughout a probability and threshold cell. A separate quartic
for each binomial count uses one common pair of conditional noise laws.
Signed second-moment information restricts their variances to a polygon.
A finite-sample fourth-moment bound makes the upper expectation convex
on that polygon. This note proves the reduction; acquisition of its
hypotheses and completion of a numerical cover are separate obligations.
\end{abstract}

\section{The conditional problem}
Throughout, $\Phi$ and $\phi$ are the standard normal distribution
function and density, and $\CE=(3+\sqrt{10})/(6\sqrt{2\pi})$.
Write $Y=B+U$, where $\Pp(B=1)=p$, $q=1-p$, and
$\E(U\mid B=i)=0$. The two conditional laws $U_0,U_1$ are fixed
throughout every calculation. Let $n\ge2$ be an integer and assume
$A_0,A_1\ge0$. Put
\[
 s_i=\E U_i^2,\quad x=qs_0,\quad y=ps_1,\quad
 v=pq,\quad w=v+x+y,\quad \lambda=n(x+y).
\]
Then $X=(Y-p)/\sqrt w$ is centered with variance one. Assume
\[
 -A_0\le U_0\le a_0,\qquad -a_1\le U_1\le A_1,\qquad
 0\le a_i\le r_i,\quad a_0\le p,\quad a_1\le q,
\]
\[
 0\le\lambda\le L,\qquad \rho=\E|X|^3\ge\rho_->0,\qquad
 b_-\le\beta:=\E[X|X|]\le b_+.
\]
The sign restrictions ensure $X\le0$ in class zero and $X\ge0$ in
class one. On a closed probability cell $[l,h]\subset(0,1/2)$ choose
support ceilings $A_i,r_i$ independently of $p$. For instance,
$a_0+a_1\le d$ permits $r_0=\min(d,h)$,
$r_1=\min(d,1-l)$; dropping the remaining joint restriction enlarges
the domain.

For a phase $u\in[u_-,u_+]$ and integer $k$, the target is the
positive closed-event score
\[
 \frac{\sqrt n}{\rho}
 \left\{\Pp\left(\sum_{i=1}^nY_i\le k+u\right)
 -\Phi\left(\frac{k+u-np}{\sqrt{nv+\lambda}}\right)\right\}<\CE.
 \tag{1}
\]
This is one oriented score. Both signs of the original discrepancy
require their own acquisition or the justified reflection argument.

\section{One common variance polygon}
The sign partition gives the exact identity
\[
 \beta=\frac{v(1-2p)-x+y}{v+x+y}.
 \tag{2}
\]
Indeed the positive and negative conditional square contributions are
$p(q^2+s_1)$ and $q(p^2+s_0)$. Since the denominator is positive,
the beta bounds are equivalent to
\begin{align*}
 -(1+b_+)x+(1-b_+)y&\le v[b_+-(1-2p)],\\
 (1+b_-)x-(1-b_-)y&\le v[(1-2p)-b_-].
\end{align*}
Replace each right side by a rational upper bound for its cubic on
$[l,h]$. Together with $x,y\ge0$ and $x+y\le L/n$, these halfplanes
define a fixed compact outer polygon $\mathcal V$. Every actual
common variance pair belongs to it, though not every point need be
realizable. Intersecting each pair of boundary lines and retaining
the intersections satisfying every halfplane gives all vertices.
A reader may require a nonempty two-dimensional polygon and reject
degenerate cases. Empty-domain exclusion needs its own exact proof.

\section{Finite common-moment tests}
For $j=0,\ldots,n$ choose quartics
$P_j(z)=\sum_{r=0}^4c_{jr}z^r$, with $c_{j4}\ge0$, satisfying
\[
 P_j(z)\ge1\quad(z\le0),\qquad P_j(z)\ge0\quad(z\ge0).
 \tag{3}
\]
Let $W_j$ sum $n-j$ independent $U_0$ variables and $j$ independent
$U_1$ variables. Conditional independence gives
\[
 \Pp\left(\sum Y_i\le k+u\right)
 \le\sum_{j=0}^n\pi_j(p)\E P_j(W_j+j-k-u),\qquad
 \pi_j(p)=\binom njp^jq^{n-j}.
 \tag{4}
\]
The left condition in (3) applies at zero, retaining every event-plane
atom. All expectations use the same two conditional laws.

\begin{lemma}[Support and variance reduction]
For each fixed probability and phase, the right side of (4) is bounded
by the maximum of the following evaluations: a vertex $(x,y)$ of
$\mathcal V$ and one pair
\[
 \gamma_0\in\{-A_0,0,r_0\},\qquad
 \gamma_1\in\{A_1,0,-r_1\}.
 \tag{5}
\]
For that evaluation define
\begin{align*}
 t_j&=(n-j)x/q+jy/p,\\
 M_{j0}&=1,\quad M_{j1}=0,\quad M_{j2}=t_j,\\
 M_{j3}&=(n-j)\gamma_0x/q+j\gamma_1y/p,\\
 M_{j4}&=(n-j)\gamma_0^2x/q+j\gamma_1^2y/p
             +3(1-1/n)t_j^2,\\
 E_j&=\sum_{r=0}^4c_{jr}\sum_{i=0}^r
                \binom ri(j-k-u)^{r-i}M_{ji},\qquad
 F=\sum_{j=0}^n\pi_j(p)E_j.
\end{align*}
\end{lemma}
\begin{proof}
Put $\mu_i=\E U_i^3$. Integrating the nonnegative support polynomials
$(U_0+A_0)(a_0-U_0)U_0^2$ and its reflected counterpart gives
\[
 \E U_0^4\le(a_0-A_0)\mu_0+A_0a_0s_0,\qquad
 \E U_1^4\le(A_1-a_1)\mu_1+A_1a_1s_1.
 \tag{6}
\]
Independence and zero means give
\[
 \E W_j^4=(n-j)\E U_0^4+j\E U_1^4
       +3t_j^2-3\{(n-j)s_0^2+js_1^2\}.
\]
Cauchy's inequality for the $n$ individual variances bounds the
quantity in braces below by $t_j^2/n$. All leading coefficients
and binomial weights are nonnegative, so this gives an upper bound
in (4).

For $s_0>0$, $\gamma_0=\mu_0/s_0\in[-A_0,a_0]$.
The mixture bound is affine in $\gamma_0$. At its endpoint $-A_0$,
the coefficient of $s_0$ in (6) is $A_0^2$, independent of $a_0$.
At the other endpoint it is $a_0^2$. The remaining expression is
convex quadratic in $a_0$, hence is bounded by its values at
$0,r_0$. The reflected argument treats class one. This proves
the nine common choices (5). A zero variance makes that conditional
variable vanish and requires no division.

For each choice (5), $F$ is affine in $x,y$ except for a sum of
nonnegative multiples of squares $t_j^2$. Thus it is convex on
$\mathcal V$, and its value at a convex combination of vertices
is no larger than the largest vertex value. The displayed fourth
moment is an extremal evaluation of an upper expression, not the
false claim $\E U_i^4\le\gamma_i^2s_i$ for an arbitrary actual law.
\end{proof}

\section{Continuous cells and exact acceptance}
Let $T_-$ be a rational lower bound, uniform on the cell, for
\[
 \Phi\left(\frac{k+u-np}{\sqrt{npq+\lambda}}\right)
                 +\frac{\CE\rho_-}{\sqrt n}.
 \tag{7}
\]
The Gaussian part must retain the numerator's sign. Put
$a=k+u_--nh$, $v_-=l(1-l)$, $v_+=h(1-h)$.
A valid lower normal argument is $a/\sqrt{nv_++L}$ for $a\ge0$,
and $a/\sqrt{nv_-}$ for $a<0$. Outward square-root and Gaussian
bounds then give (7). For a negative numerator, using only the
upper noise endpoint would give the wrong inequality.

\begin{proposition}[Exact cell criterion]
Suppose (3) holds globally and, for every polygon vertex and every
pair (5), the rational polynomial
\[
 R(p,u)=p^2(1-p)^2\{T_--F(p,u)\}
 \tag{8}
\]
is strictly positive on $[l,h]\times[u_-,u_+]$. Then (1) holds on
the entire cell for every actual conditional pair satisfying the
stated hypotheses.
\end{proposition}
\begin{proof}
The cleared denominator is positive. The lemma and (4) give
$\Pp(\sum Y_i\le k+u)<T_-$. Subtract the Gaussian CDF and apply
(7), followed by $\rho\ge\rho_-$. At zero noise the same expressions
are the literal polynomial expectations at the actual threshold;
no jump is discarded and no variance is divided out.
\end{proof}

The coordinate degrees in (8) are at most $n+4$ in $p$ and four in
$u$. The moments have denominators at most $p^2q^2$; the largest
degree is from $p^2q^2\pi_j(p)$ times a constant moment. For $n=5$
the bounds are $(9,4)$. After mapping the cell to $[0,1]^2$,
a polynomial $\sum r_{ij}s^it^j$ has tensor Bernstein coefficients
\[
 b_{ab}=\sum_{i\le a,\,j\le b}r_{ij}
          \frac{\binom ai}{\binom9i}
          \frac{\binom bj}{\binom4j}.
\]
Positive coefficients imply (8), since the basis is nonnegative
and sums to one. Inconclusive coefficients do not disprove positivity.
For (3), set $z=\pm t/(1-t)$ and multiply the appropriate ray
polynomial by $(1-t)^4$. Exact positive endpoints and zero Sturm
roots on $[0,1]$ prove positivity on the ray. Repeated-factor
cases may conservatively reject; unresolved roots cannot accept.
Floating-point optimization only proposes coefficients. Exact
acceptance reconstructs these polynomials and the complete cover.


\subsection{Keeping the Gaussian term at the same noise level}
The uniform target in (7) can lose useful information by taking its
minimum at a different noise level from the largest quartic expectation.
For a negative normal numerator there is a direct remedy.

\begin{proposition}[Concave Gaussian correction]
Suppose throughout a cell
\[
 a_{\max}=k+u_+-nl\le0,\qquad
 a_{\min}^2=(k+u_--nh)^2\le3nl(1-l).
 \tag{9}
\]
Then the cell criterion remains valid if its target is chosen separately
for each polygon vertex $(x,y)$, using a lower bound for
\[
 \Phi\left(\frac{k+u-np}{\sqrt{npq+n(x+y)}}\right)
                         +\frac{\CE\rho_-}{\sqrt n}.
 \tag{10}
\]
A rational lower Gaussian value can use numerator $a_{\min}$ and
variance $nl(1-l)+n(x+y)$. It is constant on the probability/phase cell
for this fixed vertex, so the same polynomial degree bound applies.
\end{proposition}
\begin{proof}
For fixed $p,u$ set $a=k+u-np$, $D=npq+\lambda$, and
$g(\lambda)=\Phi(a/\sqrt D)$. Direct differentiation gives
\[
 g'(\lambda)=-\frac{a\phi(a/\sqrt D)}{2D^{3/2}},\qquad
 g''(\lambda)=\frac{a\phi(a/\sqrt D)}{4D^{5/2}}
                         \left(3-\frac{a^2}{D}\right).
\]
Condition (9) implies $a\le0$ and $a^2\le3npq$, so $g$ is concave
on the entire noise interval, including the constant case $a=0$.
For each common skew choice, $F-g(n(x+y))$ is therefore convex
on the polygon. Its maximum occurs at a vertex. Check (10) at
each vertex and apply the same argument as before. For a negative
numerator, increasing either the numerator or the denominator
increases its standardized value. This justifies the stated
uniform lower bound. No endpoint choice for the normal term is
borrowed from a different variance pair.
\end{proof}

\subsection{Keeping the same law's third absolute moment}
The lower bound $\rho\ge\rho_-$ may also discard useful information.
The sign partition gives an exact expression for the third absolute
moment before normalization:
\[
 Q:=\E|Y-p|^3
   =v(1-2v)+3(px+qy)-\gamma_0x+\gamma_1y,
 \qquad \rho=\frac{Q}{(v+x+y)^{3/2}}.
 \tag{11}
\]
Here $\gamma_i s_i=\E U_i^3$; a zero variance contributes zero.
Indeed, expand $q\E(p-U_0)^3+p\E(q+U_1)^3$ and use the
zero conditional means. The actual $Q$ is nonnegative.

Fix a noise slab $[L_-,L_+]$ and a rational upper bound
$W\ge v+x+y$ on it. Let $c_0>0$ and $c_1>0$ be rational lower
bounds for $\CE/\sqrt n$ and $c_0/W^{3/2}$, respectively.
For any fixed rational $\theta\in[0,1]$,
\[
 \frac{\CE\rho}{\sqrt n}
 \ge (1-\theta)c_0\rho_-+\theta c_1Q.
 \tag{12}
\]
This follows by taking a convex combination of $\rho\ge\rho_-$
and $\rho\ge Q/W^{3/2}$. Directed square-root bounds suffice to
choose $c_0,c_1$. The sign of $Q$ justifies rounding $c_1$ down.

The cell test can therefore use the right side of (12) in place of
the constant third-moment contribution in (7) or (10). To justify
the same finite reductions, first subtract $\theta c_1Q$ from
the upper expectation on actual laws. For fixed variances this
expression is affine in each $\gamma_i$. At an inward skew
endpoint the new contribution is affine in $a_i$, leaving the
nonnegative quadratic coefficient from the quartic majorant
unchanged. Thus the same nine endpoint choices apply. For each
choice, $Q$ is affine in $x,y$, so subtracting it preserves the
convexity used at the polygon vertices, including the Gaussian
correction under (9). After multiplication by $p^2q^2$, the new
term has degree at most eight in $p$; the bound $(9,4)$ for $n=5$
is unchanged.

An artificial endpoint evaluation of (11) need not be a moment
of a probability law and can be negative. This causes no problem:
nonnegativity was used only to derive (12) for actual laws.
The resulting polynomial upper expression is then maximized on
an enlarged domain. One must not round its coefficient again
according to a sign inferred from a fictitious law.

\subsection{Using acquired mean and sign information}
Two further restrictions already present in the selected original
contract tighten the same fixed polygon and endpoint bounds.
Let $m=\E X_+=v/\sqrt w$, where the positive part follows from
the sign partition. If $m\ge m_->0$, then
\[
 \lambda\le n\left(\frac{v^2}{m_-^2}-v\right).
 \tag{13}
\]
If $2v_->m_-^2$, the right side increases with $v$ throughout
the cell. We may replace the slab ceiling by
\[
 L_+^{\rm eff}=\min\left\{L_+,\,
          n\left(\frac{v_+^2}{m_-^2}-v_+\right)\right\},
 \qquad W=v_++L_+^{\rm eff}/n.
 \tag{14}
\]
Intersect the polygon also with
$L_-/n\le x+y\le L_+^{\rm eff}/n$.
This still covers every actual law in the original slab.

Suppose the original standardized conditional supports satisfy
$X\le-\varepsilon_0$ in class zero and
$X\ge\varepsilon_1$ in class one, with $\varepsilon_i>0$.
Since $X=(B+U-p)/\sqrt w$, these conditions imply
\[
 a_0\le p-\varepsilon_0\sqrt w,\qquad
 a_1\le q-\varepsilon_1\sqrt w.
\]
Retain also the acquired ceiling $a_0+a_1\le d$.
If $s_-$ is a rational lower bound for $\sqrt{v_-+L_-/n}$,
valid fixed inward ceilings are
\[
 r_0=\min(d,h-\varepsilon_0s_-),\qquad
 r_1=\min(d,1-l-\varepsilon_1s_-).
 \tag{15}
\]
The reader checks their nonnegativity. It rejects empty or
degenerate polygons rather than treating a failed calculation as
an exclusion. These are consequences of acquired hypotheses;
their use does not independently prove those hypotheses.

\section{The selected row and the complete-cover obligation}
The original five-summand BN probability-panel38 contract supplies
\[
 p\in[19/50,39/100],\quad \rho\in[11/10,23/20],\quad L=7687/25000,
\]
\[
 (b_-,b_+)=(-15293/500000,307221/1000000).
\]
\[
 A_0=82453/125000,\qquad A_1=703209/1000000.
\]
These are acquired inputs, not consequences of the criterion.
For (13)--(15) the same contract supplies
\[
 m_-=217391/500000,\quad d=471533/500000,\quad
 \varepsilon_0=29529/10^6,\quad\varepsilon_1=9459/200000.
\]
The original raw threshold bound $|H|\le2051019/10^6$ and
$\sqrt w\le547171/10^6$ imply, for the Bernoulli-coordinate threshold
$h_B=5p+\sqrt wH$,
\[
 \frac{777741882751}{10^{12}}\le h_B\le
 \frac{3072258117249}{10^{12}}.
\]
Thus $1/2<h_B<7/2$, and only nearest thresholds $k=1,2,3$ occur.
Both closed phase halves must be retained for each. At half-integers
both representations may be included. Mixture counts $j=0,\ldots,5$
have a different role and do not justify this threshold restriction.

A completed row certificate must cover every mapped point or exclude
it using a proved original necessary condition. A finite tree must
retain both children at each split, all roots and all shared faces.
Missing roots or children, duplicated root selections, and malformed
objects reject. The earlier
fixed-probability, fixed-phase pilot does not fulfill this obligation.
The accompanying progress records identify which continuous regions
have actually been accepted; an incomplete cover cannot justify
removing a manuscript argument.

In this setting a certificate consists of a finite subdivision tree,
six rational quartics and a rational weight $\theta$ at every leaf.
Its mathematical claim is the score inequality (1) for every actual
law in the specified acquired row. A floating-point linear program
finds candidate coefficients. A separate standard-library Python
reader recomputes the ray and tensor polynomials with exact rational
arithmetic, checks strict inequalities, and checks both closed
children at every split. A source-bound collector fixes the required
threshold roots independently of the candidate data and runs the
reader afresh on each. Thus a favorable search result or a stored
success message is insufficient. The analytic reduction above,
the original acquisition, and the exact readers are distinct parts
of the proof available for human inspection.
\end{document}

